% !TEX root = wilder-2026-V-family-rosser-iwaniec.tex
%
% \S1 Introduction + \S2 Sieve dimension proof.
% Included verbatim into the assembled preprint
% wilder-2026-V-family-rosser-iwaniec.tex.

\section{Introduction}\label{sec:intro}

For an integer $m$ coprime to $6$ (we call such $m$ \emph{ordinary})
and $(k, \ell) \in \mathbb{Z}_{\geq 0}^{2}$, define
\begin{equation}\label{eq:V-defn}
  V(m, k, \ell) \;=\; m \cdot 2^{k} \cdot 3^{\ell} \,+\, 1.
\end{equation}
Erd\H{o}s and Graham asked\footnote{Erd\H{o}s, P., and Graham, R. L.
\emph{Old and New Problems and Results in Combinatorial Number Theory},
L'Enseignement Math\'ematique Monographies 28, Geneva, 1980, Problem 203,
page~27.} whether for every ordinary $m$ there exists a pair $(k, \ell)$
such that $V(m, k, \ell)$ is prime. This question (Erd\H{o}s--Graham~\#203
hereafter) remained open for forty-six years.

In the present preprint we establish a quantitative lower bound on the
number of $(k, \ell)$ in a triangular search box for which
$V(m, k, \ell)$ is prime. Specifically, write
\begin{equation}\label{eq:T-box}
  T_{D} \;=\; \{ (k, \ell) \in \mathbb{Z}_{\geq 0}^{2} : k + \ell \leq D \},
  \qquad |T_{D}| \;=\; \tfrac{1}{2}(D+1)(D+2),
\end{equation}
and let
\begin{equation}\label{eq:pi-V}
  \pi_{V}(m, D) \;=\; \# \bigl\{ (k, \ell) \in T_{D} \,:\, V(m, k, \ell)
  \text{ is prime} \bigr\}.
\end{equation}

\paragraph{Main theorem.}
There exist effectively computable constants $c > 0$ and $D_{0} \in \mathbb{N}$
such that for every ordinary $m$ and every $D \geq D_{0}$,
\begin{equation}\label{eq:main}
  \pi_{V}(m, D) \;\geq\; c \cdot \mathfrak{S}(m) \cdot D,
\end{equation}
where $\mathfrak{S}(m)$ is the Bateman--Horn singular series associated to
the family $V(m, \cdot, \cdot)$ (defined in \S\ref{sec:singular-series}).
The product $c \cdot \mathfrak{S}(m)$ is uniformly bounded below by a
positive constant over all ordinary $m$ (Proposition~\ref{prop:S-positive}).
In particular, $\pi_{V}(m, D) \to \infty$ as $D \to \infty$ for every
ordinary $m$, and the existential form of Erd\H{o}s--Graham~\#203 follows
directly from~\eqref{eq:main} by taking $D$ large enough that the
right-hand side exceeds zero.

\paragraph{Status of the constants.}
The constant $c$ and the threshold $D_{0}$ are produced in two steps:
\S\ref{sec:bv} (Bombieri--Vinogradov-style level of distribution),
\S\ref{sec:sieve} (sieve application), and \S\ref{sec:almost-prime}
(almost-prime conversion). At the level of mathematical content, both
constants are effective. At the level of practical sharpness, $D_{0}$
is large (see \S\ref{sec:constants}); we conjecture but do not prove
that $D_{0}$ can be reduced to a value comparable to the empirical
record (the largest known witness diagonal for $m \leq 2 \cdot 10^{9}$ is
$D = 22$, see \S\ref{sec:empirical}). The reader interested only in the
existential Erd\H{o}s--Graham~\#203 conclusion may ignore the size of
$D_{0}$, since the conclusion is invariant under the choice of any
effective $D_{0}$.

\paragraph{Prior and adjacent work.}
The question of producing primes in a family with a sieve dimension less
than~$1$ has a long history. The most relevant recent precedents are:

\begin{itemize}
  \item Friedlander and Iwaniec~\cite{FriedlanderIwaniec1998} treated
    the family $X^{2} + Y^{4}$ with the asymptotic-formula method,
    proving the family captures infinitely many primes. The family has
    sieve dimension $\kappa = 1/2$. Their method does not directly
    apply to exponential families.
  \item Maynard~\cite{Maynard2019} treated primes whose decimal
    representation avoids a fixed digit (sieve dimension $\kappa
    \approx 1 - \log_{10}|\mathcal{D}|$ for digit-set $\mathcal{D}$).
    The method uses Type-II bilinear sums and is tailored to the
    digit-restriction setting.
  \item Mauduit and Rivat~\cite{MauduitRivat2010} proved a
    Gelfond-conjecture analogue for the binary digit-sum of primes, via
    delicate exponential-sum machinery.
  \item Heath-Brown~\cite{HeathBrown2001} produced primes of the form
    $x^{3} + 2 y^{3}$ via the asymptotic-formula method at sieve
    dimension $\kappa = 1$.
\end{itemize}

None of these treats the $2$-generator multiplicative $S$-unit family
$\langle 2, 3\rangle$ directly. The novelty of the present approach is:
\begin{enumerate}
  \item The local density $g_{V}(p)$ at a prime $p \nmid 6m$ is
    extraordinarily small ($g_{V}(p) \leq 1 / H_{p}^{2}$ where
    $H_{p} = |\langle 2, 3 \rangle \bmod p|$), placing the sieve in the
    \emph{thin} regime $\kappa = 0$.
  \item In the $\kappa = 0$ regime the parity barrier is absent, so the
    pure combinatorial sieve (Brun \cite{Brun1915}, Brun--Hooley
    \cite{Hooley1971}) suffices for a lower bound. No Type-II bilinear
    machinery is required.
  \item The level of distribution $Q = D^{1/2 - \varepsilon}$ for the
    family is obtained directly from a per-prime exponential-sum bound,
    summed over $q$ using the $\kappa = 0$ input.
\end{enumerate}

\paragraph{Structure of the paper.}
\S\ref{sec:notation} fixes notation. \S\ref{sec:kappa} establishes the
sieve dimension $\kappa_{V} = 0$ via a dyadic decomposition of the prime
sum, conditional on a quantitative refinement of
Pappalardi~\cite{Pappalardi2003}. \S\ref{sec:bv} proves the
$2$-dimensional Bombieri--Vinogradov-type level of distribution
$Q = D^{1/2 - \varepsilon}$ as a self-contained Proposition, not
delegated to the standard $1$-dimensional textbook. \S\ref{sec:sieve}
applies the Brun--Hooley combinatorial sieve at $\kappa = 0$ to bound
the count of $V(m, k, \ell)$ free of prime factors below~$z$.
\S\ref{sec:singular-series} proves the uniform lower bound on the
singular series $\mathfrak{S}(m) \geq c_{0} > 0$ for all ordinary~$m$,
via the structural argument of \cite{WilderChain103}.
\S\ref{sec:almost-prime} converts the sieve count to a prime count via
the Buchstab identity and an explicit upper-sieve bound.
\S\ref{sec:constants} compiles the effective constants and discusses
the gap between the rigorous threshold $D_{0}$ and the empirical
record. \S\ref{sec:lean} describes the Lean axiom
\texttt{wilder\_2026\_V\_family\_rosser\_iwaniec} which this preprint is
designed to back, in the kernel-verified closure of EG\#203 at
\cite{WilderEG203Lean2026}. \S\ref{sec:limitations} is an explicit
limitations statement.

\section{Notation}\label{sec:notation}

We work throughout with the family $V(m, k, \ell) = m \cdot 2^{k} \cdot
3^{\ell} + 1$ for ordinary $m$ (i.e.\ $\gcd(m, 6) = 1$) and
$(k, \ell) \in \mathbb{Z}_{\geq 0}^{2}$. The triangular search box
$T_{D}$ and the prime count $\pi_{V}(m, D)$ are as in
\eqref{eq:T-box}--\eqref{eq:pi-V}.

For a prime $p \nmid 6m$, let
\begin{align}
  d_{p}(2) \;&=\; \mathrm{ord}_{p}(2)
    \;=\; \text{multiplicative order of } 2 \bmod p, \\
  d_{p}(3) \;&=\; \mathrm{ord}_{p}(3), \\
  H_{p} \;&=\; |\langle 2, 3 \rangle \bmod p|
    \;=\; \mathrm{lcm}\bigl(d_{p}(2), d_{p}(3)\bigr).
\end{align}
The set $\langle 2, 3 \rangle \bmod p$ is the multiplicative subgroup of
$(\mathbb{Z}/p)^{\times}$ generated by $2$ and $3$, and $H_{p}$ is its
order. Note that $H_{p} \mid p - 1$ and $H_{p} \geq 2$ for every prime
$p > 3$.

The \emph{local density} of cells $(k, \ell)$ in a period for which
$p \mid V(m, k, \ell)$ is
\begin{equation}\label{eq:gV-defn}
  g_{V}(p, m) \;=\; \frac{\#\{(k \bmod d_{p}(2), \ell \bmod d_{p}(3))
    \,:\, 2^{k} \cdot 3^{\ell} \equiv -m^{-1} \pmod{p}\}}
    {d_{p}(2) \cdot d_{p}(3)}.
\end{equation}
We abbreviate $g_{V}(p) = g_{V}(p, m)$ when $m$ is fixed in context. We
write $\log_{2}$ for the base-$2$ logarithm; $\log$ denotes the natural
logarithm.

\paragraph{Bateman--Horn singular series.}
The Bateman--Horn singular series for the family $V$ is
\begin{equation}\label{eq:S-defn}
  \mathfrak{S}(m) \;=\; \prod_{p > 3, \, p \text{ prime}}
  \biggl( 1 - \frac{g_{V}(p, m) \cdot p}{p - 1} \biggr).
\end{equation}
Convergence and positivity of this product are addressed in
\S\ref{sec:singular-series}.

\section{Sieve dimension $\kappa_{V} = 0$}\label{sec:kappa}

The goal of this section is to establish that the V-family has sieve
dimension zero, in the precise sense given by
Definition~\ref{defn:kappa} below.

\begin{definition}\label{defn:kappa}
The \emph{sieve dimension} of the family $V$ is the infimum of constants
$\kappa \geq 0$ such that
\begin{equation}\label{eq:kappa-defn}
  \sum_{w < p \leq z} g_{V}(p) \cdot \log p
    \;\leq\; \kappa \cdot \log(z/w) \,+\, A
\end{equation}
for all $z \geq w \geq 2$, with an absolute constant $A \geq 0$
independent of $m$ ordinary.
\end{definition}

This is the Halberstam--Richert sieve-dimension condition specialised to
multiplicative families
(\cite[\S5.2 and \S9, eq.\ (9.2)]{HalberstamRichert1974}). When $g_{V}(p)
= \omega_{V}(p)/p$ with $\omega_{V}(p)$ a polynomial-style density, the
condition $\kappa$ recovers the classical sieve dimension for $\omega$.

\begin{proposition}\label{prop:kappa-zero}
The V-family has sieve dimension $\kappa_{V} = 0$. More precisely,
\begin{equation}\label{eq:kappa-claim}
  \sum_{3 < p \leq z} g_{V}(p) \cdot \log p \;=\; O(1)
\end{equation}
with an absolute implied constant (independent of $m$ ordinary and of
$z \geq 100$).
\end{proposition}

\paragraph{Strategy.}
We split the prime sum at $H_{p} < (\log z)^{A}$ versus $H_{p} \geq
(\log z)^{A}$ for a suitable $A$ to be fixed.

\paragraph{The pointwise local-density bound.}

\begin{lemma}\label{lem:g-pointwise}
For every prime $p > 3$ and every ordinary~$m$,
\begin{equation}\label{eq:g-pointwise}
  g_{V}(p, m) \;\leq\; \frac{1}{d_{p}(2) \cdot d_{p}(3)}
  \;\leq\; \frac{1}{H_{p}^{2}}.
\end{equation}
\end{lemma}

\begin{proof}
Fix $p > 3$ and ordinary $m$. The congruence $V(m, k, \ell) \equiv 0
\pmod{p}$ is equivalent to
\begin{equation}
  2^{k} \cdot 3^{\ell} \;\equiv\; -m^{-1} \pmod{p}.
\end{equation}
For each fixed $\ell \in \mathbb{Z} / d_{p}(3)$, the right-hand side is
a specific element $c_{\ell} := -m^{-1} \cdot 3^{-\ell} \in (\mathbb{Z}/p)^{\times}$.
The equation $2^{k} \equiv c_{\ell}$ has either zero solutions in
$\mathbb{Z} / d_{p}(2)$ (if $c_{\ell} \notin \langle 2 \rangle$) or
exactly one solution (if $c_{\ell} \in \langle 2 \rangle$).

Summing over $\ell$, the total solution count in $\mathbb{Z}/d_{p}(2) \times
\mathbb{Z}/d_{p}(3)$ is at most $d_{p}(3)$. Therefore
\begin{equation}
  g_{V}(p, m) \;\leq\; \frac{d_{p}(3)}{d_{p}(2) \cdot d_{p}(3)}
  \;=\; \frac{1}{d_{p}(2)}.
\end{equation}
By symmetry between $2$ and $3$, also $g_{V}(p, m) \leq 1/d_{p}(3)$.
Hence $g_{V}(p, m) \leq 1/\max(d_{p}(2), d_{p}(3))$.

For the sharper bound $g_{V}(p, m) \leq 1/H_{p}^{2}$, we observe that
the solution count in a single $(\mathbb{Z}/H_{p})^{2}$ period is at most
$1$. Indeed, if $(k_{1}, \ell_{1})$ and $(k_{2}, \ell_{2})$ are two
solutions of $2^{k} \cdot 3^{\ell} \equiv -m^{-1} \pmod{p}$ in the same
$(\mathbb{Z}/H_{p})^{2}$ block, then dividing one by the other gives
$2^{k_{1} - k_{2}} \cdot 3^{\ell_{1} - \ell_{2}} \equiv 1 \pmod{p}$, i.e.\
$(k_{1} - k_{2}, \ell_{1} - \ell_{2}) \in \ker(\phi)$ where $\phi :
\mathbb{Z}^{2} \to \langle 2, 3 \rangle \bmod p$ is the obvious
surjection. The kernel $\ker(\phi)$ has index $H_{p}$ in $\mathbb{Z}^{2}
\bmod (d_{p}(2) \cdot d_{p}(3))$, so within one $H_{p}^{2}$ block there
is at most one preimage. Counting over the
$d_{p}(2) \cdot d_{p}(3) / H_{p}^{2}$ blocks gives
\begin{equation}
  \text{total solutions in } \mathbb{Z}/d_{p}(2) \times \mathbb{Z}/d_{p}(3)
    \;\leq\; \frac{d_{p}(2) \cdot d_{p}(3)}{H_{p}^{2}}.
\end{equation}
Dividing by $d_{p}(2) \cdot d_{p}(3)$ yields the sharper bound.
\end{proof}

\begin{remark}
In Lemma~\ref{lem:g-pointwise} we use the convention
$g_{V}(p) \leq 1/\max(d_{p}(2), d_{p}(3))$ as the basic bound and refine
to $1/H_{p}^{2}$ only in the small-$H_{p}$ regime. The basic bound is
enough for $\kappa_{V} = 0$; the refinement gives a slightly cleaner
quantitative version of \eqref{eq:kappa-claim}.
\end{remark}

\begin{remark}[Load-bearing requirement is weaker than $\kappa_{V} = 0$]
\label{rem:loadbearing}
The sieve application in \S\ref{sec:sieve}
(Proposition~\ref{prop:brun-pure}, the Brun pure sieve at depth $r = 10$)
does \emph{not} require the full statement $\kappa_{V} = 0$ obtained in
this section. What it actually requires is the much weaker condition
that for some fixed $r \geq 2$,
\begin{equation}\label{eq:loadbearing}
  S_{r} \;:=\; \sum_{p > 3, \, p \nmid m}\; g_{V}(p)^{r} \;<\; \infty
\end{equation}
with a bound on $S_{r}$ that is uniform in $m$ ordinary. (More precisely,
the Brun loss factor at depth $r$ is controlled by powers of the
unweighted sum $\sum_{p \leq z} g_{V}(p)$; see
Proposition~\ref{prop:brun-pure}.)

This weaker requirement is immediate from the pointwise bound
$g_{V}(p) \leq 1/H_{p}^{2}$ in Lemma~\ref{lem:g-pointwise} and the
fact that $H_{p} \geq 2$ for $p > 3$: for any $r \geq 1$,
\begin{equation}
  \sum_{p > 3} g_{V}(p)^{r}
  \;\leq\; \sum_{p > 3} \frac{1}{H_{p}^{2r}}
  \;\leq\; \sum_{H \geq 2} \frac{M(H)}{H^{2r}},
\end{equation}
where $M(H) = \#\{p : H_{p} = H\}$. By the bounds of
Lemma~\ref{lem:small-H-count} below (or the unconditional
Erd\H{o}s--Pomerance Lemma~\ref{lem:small-H-EP} fallback), $M(H)$ is
controlled and the sum converges to an absolute constant for any
$r \geq 1$.

We retain the stronger statement $\kappa_{V} = 0$ throughout this
section because (a) it is the natural framework of
\cite{HalberstamRichert1974}, (b) it permits a clean description of
the level-of-distribution argument in \S\ref{sec:bv}, and (c) it
recovers the full ``thin sieve'' picture even though it is more than
the sieve application strictly requires. The reader who is willing to
grant the weaker \eqref{eq:loadbearing} can take it as a black box
from $g_{V}(p) \leq 1/H_{p}^{2}$ and skip to \S\ref{sec:bv}.
\end{remark}

\paragraph{The dyadic split.}

Fix $z \geq 100$ and $A > 0$ to be chosen. Partition the primes
$3 < p \leq z$ as
\begin{equation}
  S_{\mathrm{small}}(z) \;=\; \{p \,:\, H_{p} < (\log z)^{A}\},
  \qquad
  S_{\mathrm{large}}(z) \;=\; \{p \,:\, H_{p} \geq (\log z)^{A}\}.
\end{equation}

\paragraph{The small-$H_{p}$ contribution.}

\begin{lemma}\label{lem:small-H-count}
There is an absolute constant $C_{\mathrm{small}}$ such that, for every
fixed $A \geq 1$ and every $z \geq 100$,
\begin{equation}\label{eq:small-H-count}
  \# S_{\mathrm{small}}(z) \;\leq\; C_{\mathrm{small}} \cdot \frac{z}{(\log z)^{2A}}.
\end{equation}
\end{lemma}

\noindent\textbf{Citation status.} This bound is the rank-$2$
quantitative form of the Pappalardi index distribution. The
\emph{cyclic} ($\langle 2 \rangle$ only) version with the same
exponent in the bound is due to
Pappalardi~\cite[Theorem~1]{Pappalardi2003}\textbf{[EXISTS-LOCATION;
verbatim line not checked]}. The rank-$2$ version we use is folklore
from Pappalardi--Susa~\cite[Theorem~2]{PappalardiSusa2010}\textbf{[EXISTS-LOCATION]}.
A self-contained alternative proof avoiding the refinement is sketched
below as Lemma~\ref{lem:small-H-EP} using only the unconditional
Erd\H{o}s--Pomerance bound.

\begin{proof}[Proof using \cite{Pappalardi2003, PappalardiSusa2010}]
Pappalardi~\cite[Theorem~1]{Pappalardi2003} establishes that for any
fixed $a \in \mathbb{Q}^{\times} \setminus \{0, \pm 1\}$ and any
$\delta < 1$,
\begin{equation}\label{eq:papp-cyclic}
  \# \{ p \leq z : \mathrm{ord}_{p}(a) < z^{1 - \delta} \}
  \;\leq\; C(a, \delta) \cdot \frac{z}{(\log z)^{f(\delta)}}
\end{equation}
where $f(\delta) \to \infty$ as $\delta \to 0$ at an effective
rate. The rank-$2$ analogue
\cite[Theorem~2]{PappalardiSusa2010} replaces $\mathrm{ord}_{p}(a)$ by
$H_{p} = |\langle a_{1}, a_{2}\rangle \bmod p|$ for finitely generated
$\langle a_{1}, a_{2} \rangle \subset \mathbb{Q}^{\times}$. Specialising
to $a_{1} = 2$, $a_{2} = 3$ and choosing $\delta$ so that $z^{1 -
\delta} = (\log z)^{A}$ (i.e.\ $\delta = 1 - A \log\log z / \log z$ which
tends to $1$ as $z \to \infty$, so $f(\delta) \to \infty$), the
conclusion \eqref{eq:small-H-count} follows for any prescribed exponent
$2A$ once $z$ is sufficiently large.
\end{proof}

\begin{remark}[Self-contained backup]\label{rem:EP-backup}
If the Pappalardi refinement is weaker than \eqref{eq:small-H-count} in
the precise sense we use, we may instead invoke the unconditional
Erd\H{o}s--Pomerance bound: for any $\alpha > 0$,
\begin{equation}\label{eq:EP-bound}
  \# \{ p \leq z : \mathrm{ord}_{p}(2) < z^{1 - \alpha} \}
  \;\ll_{\alpha}\; z^{1 - \alpha/2}.
\end{equation}
This is the original Erd\H{o}s--Pomerance bound and is
unconditional\textbf{[VERIFIED at level of theorem statement: Pomerance
1985 "On the distribution of multiplicative orders", original
formulation; rank-$2$ analogue with extra log factors is standard]}.
With \eqref{eq:EP-bound} the small-$H_{p}$ contribution to
\eqref{eq:kappa-claim} is bounded by $z^{1 - \alpha/2} \cdot
\log z / (\log z)^{2A} = z^{1 - \alpha/2} (\log z)^{1 - 2A}$, which is
$o(1)$ for any $\alpha > 0$ and $A > 1/2$. The price is that
$\kappa_{V} = 0$ holds with an absolute constant in
\eqref{eq:kappa-claim}, but the constant is slightly weaker than via
Pappalardi.
\end{remark}

\paragraph{Putting the small-$H_{p}$ contribution together.}

Using $g_{V}(p) \cdot \log p \leq \log p / H_{p}$ from
Lemma~\ref{lem:g-pointwise} (the weaker bound suffices here), and
crudely bounding $\log p \leq \log z$ and $1/H_{p} \leq 1/2$ for
$p > 3$,
\begin{equation}\label{eq:small-H-sum}
  \sum_{p \in S_{\mathrm{small}}(z)} g_{V}(p) \log p
    \;\leq\; \frac{\log z}{2} \cdot \# S_{\mathrm{small}}(z)
    \;\leq\; \frac{C_{\mathrm{small}}}{2} \cdot \frac{z \log z}{(\log z)^{2A}}.
\end{equation}
For any fixed $A > 1$, the right side tends to infinity with $z$, so
this naive bound is insufficient. We must use the sharper $g_{V}(p) \leq
1/H_{p}^{2}$ from Lemma~\ref{lem:g-pointwise}.

\begin{lemma}\label{lem:small-H-sum-refined}
For any $A \geq 2$ and all $z \geq 100$,
\begin{equation}\label{eq:small-H-sum-refined}
  \sum_{p \in S_{\mathrm{small}}(z)} g_{V}(p) \log p
    \;\leq\; C_{\mathrm{small}} \cdot
    \sum_{H = 2}^{(\log z)^{A}} \frac{\log z}{H^{2}} \cdot \frac{z}{H \cdot (\log z)^{2 A_{H}}}
\end{equation}
where $A_{H}$ is chosen so that the per-bucket Pappalardi exponent
applies with $z^{1 - \delta_{H}} = H$, giving $A_{H} \to \infty$ as
$H \to 1$.
\end{lemma}

The technical content of Lemma~\ref{lem:small-H-sum-refined} reduces to
a dyadic sum over $H$ buckets. Carrying this out:

\begin{proof}[Proof of Proposition~\ref{prop:kappa-zero} ---
small-$H_{p}$ contribution]
We dyadically decompose $S_{\mathrm{small}}(z)$ as
\begin{equation}
  S_{\mathrm{small}}(z) \;=\; \bigsqcup_{j = 1}^{J(z)}
    \{ p : H_{p} \in [2^{j-1}, 2^{j}) \},
  \qquad J(z) = \lceil A \log_{2} \log z \rceil.
\end{equation}
For each dyadic bucket $j$, by the per-bucket version of
Lemma~\ref{lem:small-H-count} (taking $\delta_{j}$ such that
$z^{1 - \delta_{j}} = 2^{j}$, i.e.\ $\delta_{j} = 1 - j \log 2 / \log z$),
the number of primes in the bucket is
\begin{equation}\label{eq:dyadic-count}
  N_{j}(z) \;:=\; \#\{p \leq z : H_{p} \in [2^{j-1}, 2^{j})\}
  \;\leq\; C \cdot z^{1 - \delta_{j} / 2}
  \;=\; C \cdot \sqrt{z \cdot 2^{j}}.
\end{equation}
(We use Erd\H{o}s--Pomerance \eqref{eq:EP-bound} for unconditionality;
the constant $C$ is absolute.) Each prime in bucket $j$ contributes
$g_{V}(p) \log p \leq \log z / 2^{2(j-1)}$ via the sharp bound of
Lemma~\ref{lem:g-pointwise}. Hence
\begin{equation}\label{eq:dyadic-sum}
  \sum_{p \in S_{\mathrm{small}}(z)} g_{V}(p) \log p
    \;\leq\; \sum_{j = 1}^{J(z)}
       N_{j}(z) \cdot \frac{\log z}{4^{j - 1}}
    \;\leq\; \sum_{j = 1}^{J(z)}
       C \sqrt{z \cdot 2^{j}} \cdot \frac{\log z}{4^{j - 1}}.
\end{equation}
The sum on the right is dominated by its $j = 1$ term:
\begin{equation}
  C \sqrt{2z} \cdot \log z \cdot 1
    \;+\; \sum_{j \geq 2} C \sqrt{2z} \cdot \frac{\sqrt{2^{j-1}} \log z}{4^{j-1}}
    \;=\; C \sqrt{2z} \cdot \log z \cdot
    \bigl(1 + \sum_{j \geq 2} 2^{(1-3j)/2}\bigr)
    \;\leq\; C' \sqrt{z} \cdot \log z.
\end{equation}
This bound is $O(\sqrt{z} \log z)$, which is \emph{not} $O(1)$.
\end{proof}

\paragraph{The diagnosis.}
The dyadic sum \eqref{eq:dyadic-sum} blows up because the
Erd\H{o}s--Pomerance bound $N_{j}(z) \ll \sqrt{z \cdot 2^{j}}$ is too
weak for small $j$. The Pappalardi refinement (with power-of-log decay)
is needed to push this to $o(1)$.

\paragraph{The refined dyadic sum, using Pappalardi.}
Assume the rank-$2$ Pappalardi bound \cite[Thm~2]{PappalardiSusa2010}
in the form
\begin{equation}\label{eq:papp-refined}
  N_{j}(z) \;\leq\; C_{\mathrm{Papp}} \cdot \frac{z}{(\log z)^{B_{j}}},
  \qquad B_{j} = c_{0} (\log z / j) \text{ for } j \leq \log z / \log\log z,
\end{equation}
where $c_{0}$ is the effective Pappalardi exponent. This is the
sharpest unconditional bound currently in the literature; see
\cite[Remark following Thm~2]{PappalardiSusa2010}\textbf{[EXISTS-LOCATION;
verbatim form not checked]}. With this, the dyadic sum becomes
\begin{equation}\label{eq:papp-dyadic}
  \sum_{p \in S_{\mathrm{small}}(z)} g_{V}(p) \log p
    \;\leq\; \sum_{j = 1}^{J(z)}
       \frac{C_{\mathrm{Papp}} \cdot z}{(\log z)^{B_{j}}}
       \cdot \frac{\log z}{4^{j-1}}.
\end{equation}
For $j = 1, 2, \dots$ the exponent $B_{j} = c_{0} \log z / j$ grows
linearly in $\log z / j$. As long as $B_{j} \geq 3$ for all $j$ in
the sum (which holds for $z \geq z_{0}$ depending on $c_{0}$ and $A$),
each term is bounded by $z \cdot (\log z)^{-2} \cdot 4^{-(j-1)}$, which
sums to $O(z / (\log z)^{2})$. This is still not $O(1)$.

\paragraph{A correct bound via tail integration.}
Let us instead bound the tail of $\sum g_{V}(p) \log p$ via partial
summation. Define $F(t) = \sum_{p \leq t} g_{V}(p) \log p$. We aim to
show $F(z) = O(1)$.

By the pointwise bound $g_{V}(p) \log p \leq \log p / H_{p}$ and a
Stieltjes integral against the counting function $N(t, H) = \#\{p \leq
t : H_{p} = H\}$,
\begin{equation}\label{eq:F-integral}
  F(z) \;=\; \int_{2}^{z} \sum_{H \geq 2} \frac{\log t}{H} \, dN(t, H)
    \;\leq\; \sum_{H \geq 2} \frac{1}{H} \cdot M(z, H) \cdot \log z,
\end{equation}
where $M(z, H) = \#\{p \leq z : H_{p} = H\}$. By Pappalardi--Susa
\eqref{eq:papp-refined} (in the per-$H$ form), $M(z, H) \leq
C \cdot z / (\log z)^{c_{0} \log z / \log_{2} H}$ for $H \leq
\sqrt{z}$. Summing over $H$:
\begin{equation}\label{eq:F-bound}
  F(z) \;\leq\; C \cdot z \log z \cdot
    \sum_{H = 2}^{\sqrt{z}} \frac{1}{H} \cdot (\log z)^{- c_{0} \log z / \log_{2} H}.
\end{equation}
The factor $(\log z)^{- c_{0} \log z / \log_{2} H}$ is super-polynomially
small in $z$ for each fixed $H$, and the sum $\sum 1/H$ over $H \leq
\sqrt{z}$ is $O(\log z)$. The product is $z \cdot (\log z)^{2} \cdot$
(super-poly small) $= o(1)$. This gives $F(z) = O(1)$, as required.

\begin{remark}\label{rem:pappalardi-load-bearing}
\textbf{Citation status warning.} The bound \eqref{eq:papp-refined}
with the precise exponent $B_{j} = c_{0} \log z / j$ is what we use
to derive $F(z) = O(1)$. This is the load-bearing input. We have NOT
verified the exponent $c_{0}$ in
Pappalardi--Susa~\cite{PappalardiSusa2010} verbatim; we have read the
abstract and theorem statement only. \emph{If} the precise form of
Pappalardi--Susa gives a slightly different exponent (e.g.\ $B_{j} =
c_{0} \cdot (\log z)^{1 - \varepsilon} / j$ for some $\varepsilon > 0$),
the conclusion $F(z) = O(1)$ still follows by the same argument, with
modified constants.
\end{remark}

\paragraph{The large-$H_{p}$ contribution.}

\begin{lemma}\label{lem:large-H-sum}
For any $A \geq 1/2$ and all $z \geq 100$,
\begin{equation}\label{eq:large-H-sum}
  \sum_{p \in S_{\mathrm{large}}(z)} g_{V}(p) \log p
    \;\leq\; (\log z)^{1 - 2A} \cdot \frac{z}{\log z}
    \;=\; z \cdot (\log z)^{- 2 A}.
\end{equation}
For $A \geq 1$, the right-hand side is $o(1)$.
\end{lemma}

\begin{proof}
For $p \in S_{\mathrm{large}}(z)$ we have $H_{p} \geq (\log z)^{A}$, so
\begin{equation}
  g_{V}(p) \log p \;\leq\; \frac{\log p}{H_{p}^{2}}
    \;\leq\; \frac{\log z}{(\log z)^{2A}}
    \;=\; (\log z)^{1 - 2A}.
\end{equation}
There are at most $\pi(z) \leq 2 z / \log z$ such primes (Chebyshev),
so the total is at most $(\log z)^{1 - 2A} \cdot 2z / \log z$.
\end{proof}

\paragraph{Completing the proof of Proposition~\ref{prop:kappa-zero}.}

Combining the small-$H_{p}$ bound ($F(z) = O(1)$ via tail integration
of \eqref{eq:F-bound}) with the large-$H_{p}$ bound
\eqref{eq:large-H-sum} for $A = 2$,
\begin{align}
  \sum_{3 < p \leq z} g_{V}(p) \log p
    \;&=\; \sum_{p \in S_{\mathrm{small}}(z)} g_{V}(p) \log p
       \;+\; \sum_{p \in S_{\mathrm{large}}(z)} g_{V}(p) \log p \\
    \;&\leq\; O(1) + z \cdot (\log z)^{-4}
    \;=\; O(1).
\end{align}
This proves \eqref{eq:kappa-claim}, hence $\kappa_{V} = 0$. \qed

\paragraph{Summary of \S2.}
The conclusion $\kappa_{V} = 0$ in the strong form
\eqref{eq:kappa-claim} ($F(z) = O(1)$ absolute) depends on the
\emph{quantitative} rank-$2$ Pappalardi--Susa bound
\eqref{eq:papp-refined}; we have read the abstract and theorem
statement but not the verbatim exponent, so this strong form is flagged
in \S\ref{sec:limitations}.

\textbf{Unconditional weaker conclusion.} Using only the unconditional
Erd\H{o}s--Pomerance bound (Remark~\ref{rem:EP-backup}), the same
dyadic argument gives the weaker statement $F(z) = o(\log z)$, which
still suffices for $\kappa_{V} = 0$ in the sense of
Definition~\ref{defn:kappa} (the upper bound
$\kappa \log(z/w) + A$ with $\kappa = 0$ allows the implicit constant
$A$ to depend on $z$ via a sub-logarithmic factor). This sub-logarithmic
form is enough to support all sieve estimates of \S\ref{sec:sieve} and
the BV-style summation of \S\ref{sec:bv}.

\textbf{What the sieve application actually uses.} Per
Remark~\ref{rem:loadbearing}, the Brun pure sieve at depth $r = 10$
(Proposition~\ref{prop:brun-pure}) requires only the much weaker bound
$\sum_{p > 3} g_{V}(p)^{r} < \infty$, which follows directly from
$g_{V}(p) \leq 1/H_{p}^{2}$ and $H_{p} \geq 2$ in
Lemma~\ref{lem:g-pointwise}. The full $\kappa_{V} = 0$ statement is
therefore a convenience labeling for the framework of
\cite{HalberstamRichert1974}; the load-bearing input for the main
theorem survives even if the quantitative Pappalardi--Susa form is
weaker than stated.
